\begin{figure*}[htbp]
    \centering
    \small
    \hspace{-0.5cm}%
    \begin{tikzpicture}[
        node distance=1.2cm,
        every node/.style={
            rectangle, 
            rounded corners,
            draw, 
            minimum width=2.5cm, 
            minimum height=0.6cm, 
            font=\footnotesize,
            align=center,
            fill=white
        },
        arrow/.style={-{Stealth[length=2mm]}},
        darrow/.style={
            -{Stealth[length=3mm]},
            dashed, line width=0.57mm,
        },
        dinterm/.style={
            dashed, line width=0.75mm,
        },
        scale=1.03
    ]
    
    % Background rectangles for each column
    \fill[RoyalBlue!20, rounded corners=5pt] (-3, 5.5) rectangle (3.25, -7.5);
    \fill[Peach!20, rounded corners=5pt] (3.25, 5.5) rectangle (7.75, -5);
    \fill[lime!20, rounded corners=5pt] (7.75, 5.5) rectangle (14, -5);
    \fill[Rhodamine!20, rounded corners=5pt] (3.25, -5) rectangle (14, -7.5);
    
    % Column headers
    \node[font=\bfseries\normalsize, fill=none, draw=none] (dataprep) at (-0.125, 5) {Dataset \\ preparation};
    \node[font=\bfseries\normalsize, fill=none, draw=none] (seqproc) at (5.5, 5) {Sequence \\ preprocessing};
    \node[font=\bfseries\normalsize, fill=none, draw=none] (phylo) at (10.875, 5) {Phylogenetic \\ analysis};
    \node[font=\bfseries\normalsize, fill=none, draw=none] (postproc) at (8.625, -5.5) {Post-processing};
    
    % Column 1 - Dataset preparation
    \node (raw) at (-0.125, 3.5) {Raw metadata};
    \node (aug) at (-0.125, 1.5) {Data augmentation};
    \node (download) at (-0.125, -0.5) {Download structures};
    \node (clean) at (-0.125, -4) {Cleaning \\ (remove short or \\ heteromeric structures)};
    \node (3di) at (-1.375, -6) {Extract \\ 3Di sequences};
    \node[minimum width=2cm] (aa) at (1.625, -6) {Extract \\ AA sequences};
    
    % Column 1 arrows
    \draw[arrow] (raw) -- (aug) node[midway, right, font=\small, fill=none, draw=none] {$n = 1,932$};
    \draw[arrow] (raw) -- (aug) node[midway, left, font=\scriptsize, fill=none, draw=none] {\textit{Cov3D}~\cite{gowthaman2021}};
    \draw[arrow] (aug) -- (download) node[midway, right, font=\small, fill=none, draw=none] {$n = 2,113$};
    \draw[arrow] (aug) -- (download) node[midway, left, font=\scriptsize, fill=none, draw=none] {\textit{PDB~\cite{rose2010} sequence} \\ similarity searches};
    \draw[arrow] (download) -- (clean) node[midway, left, font=\scriptsize, fill=none, draw=none] {\textit{biotite}~\cite{kunzmann2018}};
    \draw[arrow] (download) -- (clean) node[midway, right, fill=none, draw=none] {\includegraphics[width=2.1cm]{img/6VXX.png}};
    \draw[arrow] (clean) -- (3di);
    \draw[arrow] (clean) -- (aa) node[midway, right, font=\small, fill=none, draw=none] {$n = 676$};
    
    % Column 2 - Sequence preprocessing
    \node (align) at (5.5, 3.5) {Align};
    \node (trim) at (5.5, 2) {Remove gappy sites};
    \node (deduplicate) at (5.5, 0.5) {Deduplicate \\ (and drop gappy \\ sequences)};
    \node (subset) at (5.5, -1.25) {Subset @ x \\ Sample $x$ sequences \\ per virus/major variant};
    % \node (droplowgap) at (8, -1) {Drop gappy sequences};
    \node[minimum width=1.8cm] (subset1) at (4.5, -3.5) {Subset @ 1};
    \node[minimum width=1.8cm] (subset3) at (6.5, -3.5) {Subset @ 3};
    
    % Column 2 arrows
    \draw[arrow] (align) -- (trim) node[midway, left, font=\scriptsize, fill=none, draw=none] {\textit{MAFFT}~\cite{katoh2013}};
    \draw[arrow] (trim) -- (deduplicate) node[midway, left, font=\scriptsize, fill=none, draw=none] {\textit{trimal}~\cite{capella2009}};
    \draw[arrow] (deduplicate) -- (subset) node[midway, right, font=\small, fill=none, draw=none] {$n = 257$};
    \draw[arrow] (subset) -- (subset1) node[midway, left, xshift=0.4cm, font=\small, fill=none, draw=none] {$n = 36$};
    \draw[arrow] (subset) -- (subset3) node[midway, right, xshift=-0.4cm, font=\small, fill=none, draw=none] {$n = 62$};
    
    % Column 3 - Phylogenetic analysis
    \node[text width=2cm] (iqtreebase) at (12.5, 3.5) {Subset @ 3: \\determine substitution model};
    \node (cut) at (9.5, 3.5) {Subset @ 1: \\ cut sequences \\ into $m$ chunks \\ (150 residues)};
    \node (iqtreecut) at (11, 1.5) {Subset @ 1: \\ run IQ-TREE~\cite{minh2020} \\ on each chunked MSA};
    \node (maststart) at (11, -0.25) {Subset @ 1: \\ run MAST~\cite{wong2024} \\ using all $m$ trees};
    \node (mastiter) at (11, -2.25) {Subset @ 1: \\ repeat MAST with top \\ $m-1, \ldots, 2$ trees \\ w.r.t. likelihood};
    \node (bic) at (11, -4.25) {Model selection \\ (criterion: BIC)};
    
    % Column 3 arrows
    \draw[arrow] (iqtreebase) -- (iqtreecut) node[midway, right, xshift=-0.8cm, font=\tiny, fill=none, draw=none, color=black] {\textit{MFP}~\cite{kalyaanamoorthy2017}};
    \draw[arrow] (cut) -- (iqtreecut);
    \draw[arrow] (iqtreecut) -- (maststart);
    \draw[arrow] (iqtreebase) to [bend left=40] (maststart);
    \draw[arrow] (iqtreebase) to [bend left=45] (mastiter);
    \draw[arrow] (iqtreecut) to [bend right=65] (mastiter);
    \draw[arrow] (maststart) -- (mastiter);
    \draw[arrow] (mastiter) -- (bic);

    % Block 4 - postprocessing
    \node (domain) at (5.5, -6.5) {Protein domain \\ extraction};
    \node[fill=none, draw=none, font=\scriptsize] at (5.5, -7.1) {\textit{InterProScan}~\cite{jones2014}};
    \node (treefeat) at (11, -6.5) {Phylogenetic tree \\ feature extraction \\ (e.g., clock-likeness~\cite{zuckerkandl1962})};
    
    % Curved arrows between columns
    \draw[dinterm, blue!70] (aa.east) -- (3.25, -6);
    \draw[dinterm, blue!70] (3.25, -6) -- (3.25, 5);
    \draw[dinterm, blue!70] (3di.south) -- (-1.375, -7) node[midway, left, xshift=0.4cm, font=\scriptsize, fill=none, draw=none, color=black] {\textit{Foldseek}~\cite{vankempen2024}};
    \draw[dinterm, blue!70] (-1.375, -7) -- (3.25, -7);
    \draw[dinterm, blue!70] (3.25, -7) -- (3.25, -6);
    \draw[darrow, blue!70] (3.25, 5) -- (seqproc.west);
    \draw[dinterm, red!70] (subset3.east) -- (7.75, -3.5);
    \draw[dinterm, red!70] (subset1.south) -- (4.5, -4);
    \draw[dinterm, red!70] (4.5, -4) -- (7.75, -4);
    \draw[dinterm, red!70] (7.75, -4) -- (7.75, -1.5);
    \draw[dinterm, red!70] (7.75, -1.5) -- (7.75, 5);
    \draw[darrow, red!70] (7.75, 5) -- (phylo.west);
    \draw[darrow, red!70] (subset1.south) to [bend right=40] (domain.north);
    \draw[darrow, OliveGreen!80] (bic.south) to [bend left=40] (treefeat.north);
    
    \end{tikzpicture}
    \caption{Data processing and analysis workflow for the analysis of experimental spike glycoprotein structures from betacoronaviruses. The molecular structure image of PDB ID 6VXX~\cite{walls2019}, obtained from the RCSB PDB (\url{https://rcsb.org}) using Mol*~\cite{sehnal2021}, depicts a spike protein of a wild type SARS-CoV-2 bound with 2-acetamido-2-deoxy-beta-D-glucopyranose (NAG). MFP: ModelFinder Plus~\cite{kalyaanamoorthy2017}. BIC: Bayesian Information Criterion.}
    \label{fig:bcov3d_dag}
\end{figure*}